\documentclass[10pt]{article}
\usepackage[T1]{fontenc}
\usepackage[utf8]{inputenc}
\usepackage[french]{babel}
\usepackage{amsmath,amssymb,mathrsfs,mathpazo}
\usepackage{geometry}
\geometry{a4paper,margin=2cm}
\usepackage{graphicx,array,tabularx,tabularray,colortbl,multirow,multicol,booktabs,enumitem,pifont,ulem,fancybox,diagbox,cancel}
\usepackage{pstricks,pstricks-add,pst-plot,pst-tree,pst-node,pst-text,pst-func,pst-math,pst-bspline,pst-3dplot,pst-eucl,pst-solides3d}
\let\pscalculate\undefined
\usepackage{pst-bezier}
\usepackage{tikz,pgfplots,tkz-tab}
\usetikzlibrary{arrows,arrows.meta,patterns,calc,positioning,trees,decorations.markings,shapes.geometric,shapes.misc,shapes.arrows,intersections,decorations.pathmorphing,decorations.pathreplacing,backgrounds,shadows,babel}
\pgfplotsset{compat=1.18}
\usepgfplotslibrary{fillbetween}
\usepackage[np]{numprint}
\usepackage{listings,makecell,pas-tableur,twoopt}
\let\barre\undefined
\usepackage[pstricks]{tablvar}
\definecolor{darkgreen}{RGB}{0,100,0}
\providecommand{\red}{\color{red}}
\providecommand{\blue}{\color{blue}}
\providecommand{\green}{\color{green}}
\providecommand{\black}{\color{black}}
\providecommand{\white}{\color{white}}
\providecommand{\gray}{\color{gray}}
\providecommand{\cyan}{\color{cyan}}
\providecommand{\magenta}{\color{magenta}}
\providecommand{\yellow}{\color{yellow}}
\providecommand{\violet}{\color{violet}}
\providecommand{\orange}{\color{orange}}
\providecommand{\euro}{\texteuro}
\providecommand{\EUR}{\texteuro}
\providecommand{\ieme}{\textsuperscript{e}}
\providecommand{\degres}{\ensuremath{^\circ}}
\providecommand{\degree}{\ensuremath{^\circ}}
\providecommand{\vv}[1]{\overrightarrow{#1}}
\providecommand{\Coord}[1]{#1}
\providecommand{\decosix}{}
\providecommand{\starredbullet}{\ensuremath{\bullet}}
\pagestyle{empty}
\setlength{\parindent}{0pt}
\usepackage{adjustbox}\DeclareUnicodeCharacter{2236}{:}\DeclareUnicodeCharacter{2190}{\ensuremath{\leftarrow}}
\providecommand{\R}{\mathbb{R}}
\providecommand{\N}{\mathbb{N}}
\providecommand{\Z}{\mathbb{Z}}
\providecommand{\Q}{\mathbb{Q}}
\providecommand{\D}{\mathbb{D}}
\providecommand{\C}{\mathbb{C}}
\providecommand{\e}{{\mathrm{e}}}
\providecommand{\vect}[1]{\overrightarrow{#1}}
\providecommand{\vectt}[1]{\overrightarrow{\mathrm{#1}}}
\providecommand{\vvt}[1]{\overrightarrow{\mathrm{#1}}}
\providecommand{\np}[2][]{{#2~#1}}
\providecommand{\vv}[1]{{\overrightarrow{#1}}}
\renewcommand{\d}{{\mathrm{d}}}
\providecommand{\degres}{{^\circ}}
\providecommand{\degre}{{^\circ}}
\providecommand{\blue}{\color{blue}}
\providecommand{\red}{\color{red}}
\providecommand{\enskip}{\quad}
\providecommand{\llbracket}{\lbrack\!\lbrack}
\providecommand{\rrbracket}{\rbrack\!\rbrack}
\providecommand{\euro}{{\text{€}}}
\providecommand{\no}{{\text{n°}}}
\providecommand{\interval}[2]{{[#1;#2]}}
\providecommand{\intervFO}[2]{{[#1;#2[}}
\providecommand{\dotfill}{\ldots}

\providecommand{\ds}{\displaystyle}
\providecommand{\textding}[1]{#1}
\providecommand{\Oij}{\(\mathrm{O};\vec{\imath},\vec{\jmath}\)}
\providecommand{\Oijk}{\(\mathrm{O};\vec{\imath},\vec{\jmath},\vec{k}\)}
\providecommand{\Ouv}{\(\mathrm{O};\vec{u},\vec{v}\)}
\providecommand{\intervFF}[2]{\ensuremath{[#1;#2]}}
\renewcommand{\euro}{\ifmmode\text{\texteuro}\else\texteuro\fi}
\DeclareUnicodeCharacter{00B0}{\ifmmode^{\circ}\else\textdegree\fi}
\DeclareUnicodeCharacter{20AC}{\ifmmode\text{\texteuro}\else\texteuro\fi}
% Énoncés : Éducation nationale. Transcription et corrigés : APMEP, Denis Vergès et contributeurs. Découpage et classement : CoopMaths, Sébastien Lozano, Éric Elter et contributeurs. Adaptation de la présentation : Maths Explorer.
% CC BY-SA : https://coopmaths.fr/www/about/
% https://coopmaths.fr/alea/static/bac/2022/tex/bac_2022_05_sujet2_madagascar_4.tex
\begin{document}
\noindent\textit{APMEP / CoopMaths. CC BY-SA. Présentation adaptée par Maths Explorer.}\par\medskip
\medskip

\textbf{Partie A}

\medskip

On considère la fonction $f$ définie pour tout réel $x$ de ]0~;~1] par:

\[f(x) = \text{e}^{-x} + \ln (x).\]

\smallskip

\begin{enumerate}
\item Calculer la limite de $f$ en $0$.
\item On admet que $f$ est dérivable sur ]0~;~1]. On note $f'$ sa fonction dérivée.

Démontrer que, pour tout réel $x$ appartenant à ]0~;~1], on a :

\[f'(x) = \dfrac{1 - x\text{e}^{-x}}{x}\]

\item Justifier que, pour tout réel $x$ appartenant à ]0~;~1], on a $x\text{e}^{-x} < 1$.

En déduire le tableau de variation de $f$ sur ]0~;~1].
\item Démontrer qu'il existe un unique réel $\ell$ appartenant à ]0~;~1] tel que $f(\ell) = 0$.
\end{enumerate}

\medskip

\textbf{Partie B}

\medskip

\begin{enumerate}
\item On définit deux suites $\left(a_n\right)$ et $\left(b_n\right)$ par:

\[\left\{\begin{array}{l c l}
a_0& =& \dfrac{1}{10}\\
b_0& =& 1
\end{array}\right.\: \text{et, pour tout entier naturel }\: n,\left\{\begin{array}{l c l}
a_{n+1}&=&\text{e}^{-b_n}\\
b_{n+1}&=&\text{e}^{-a_n}
\end{array}\right.\]

	\begin{enumerate}
		\item Calculer $a_1$ et $b_1$. On donnera des valeurs approchées à $10^{-2}$ près.
		\item On considère ci-dessous la fonction \texttt{termes}, écrite en langage Python.
		
\begin{center}

\begin{tabular}{|l|}\hline
\texttt{def termes (n) :}\\
\quad \texttt{a=1/10}\\
\quad \texttt{b=1}\\
\quad \texttt{for k in range(0,n) :}\\
\qquad \texttt{c= \ldots}\\
\qquad \texttt{b  = \ldots}\\
\qquad \texttt{a = c}\\
\quad \texttt{return(a,b)}\\ \hline
\end{tabular}
\end{center}

Recopier et compléter sans justifier le cadre ci-dessus de telle sorte que la fonction termes calcule les termes des suites $\left(a_n\right)$ et $\left(b_n\right)$.
	\end{enumerate}
\item On rappelle que la fonction $x \longmapsto \text{e}^{-x}$ est décroissante sur $\R$.
	\begin{enumerate}
		\item Démontrer par récurrence que, pour tout entier naturel $n$, on a :
		
		\[0 < a_n \leqslant a_{n+1} \leqslant b_{n+1} \leqslant b_n \leqslant 1\]
		
		\item En déduire que les suites $\left(a_n\right)$ et $\left(b_n\right)$ sont convergentes.
	\end{enumerate}	
\item On note $A$ la limite de $\left(a_n\right)$ et $B$ la limite de $\left(b_n\right)$.

On admet que $A$ et $B$ appartiennent à l'intervalle ]0~;~1], et que $A = \text{e}^{-B}$ et $B = \text{e}^{-A}$.
	\begin{enumerate}
		\item Démontrer que $f(A) = 0$.
		\item Déterminer $A - B$.
	\end{enumerate}
\end{enumerate}

\end{document}
